\centerline{\bf \Huge Арифметические задачи}
\title{Арифметические задачи}


1. Петя в три раза старше Ани, а Аня на 8 лет младше Пети. Опре- делите, сколько лет каждому. Ответ обоснуйте.

2. У Маши есть 4 куска пластилина красного цвета, 3 куска пла- стилина синего цвета и 5 кусков пластилина жёлтого цвета. Сначала она разделила пополам каждый не красный кусок пластилина, а затем разделила пополам каждый не жёлтый кусок пластилина. Сколько кусков пластилина получила Маша?

3. Четыре девочки поют песни, аккомпанируя друг другу. Каждый раз одна из них играет на фортепиано, а остальные три поют. Вечером они посчитали, что Аня спела 8 песен, Таня — 6 песен, Оля — 3 песни, а Катя — 7 песен. Сколько раз аккомпанировала Таня? Обоснуйте свой ответ.

4. Часы Безумного Шляпника спешат на 15 минут в час, а часы Мартовского Зайца отстают на 10 минут в час. Однажды они поставили свои часы по часам Сони (которые остановились и всегда показывают 12:00) и договорились собраться в 5 часов вечера на традиционный файв-о-клок. Сколько времени Безумный Шляпник будет ждать Мартовского Зайца, если каждый приходит ровно в 17:00 по своим часам?

5. Два маляра красят 15-метровый коридор. Каждый из них движется от начала коридора в его конец и в какой-то момент начинает красить, пока не кончится краска. У первого маляра есть красная краска, и её хватит на покраску 9 метров коридора; у второго — жёлтая краска, и её хватит на 10 метров.
Первый маляр начинает красить, когда оказывается в двух метрах от начала коридора; а второй заканчивает красить, когда находится в метре от конца коридора. Сколько метров коридора покрашено ровно в один слой?

6. В двух кошельках лежат две монеты, причём в одном кошельке монет вдвое больше, чем в другом. Как такое может быть?

7. Надя испекла пирожки с малиной, черникой и клубникой. Пирож- ков с малиной получилась половина от общего количества пирожков; пирожков с черникой — на 14 меньше, чем пирожков с малиной. А пирожков с клубникой получилось в два раза меньше, чем пирожков с малиной и черникой вместе. Сколько пирожков каждого вида испекла Надя?

8. Аня, Боря, Вика и Гена дежурят в течение 20 дней в школе. Известно, что каждый день дежурят ровно трое из них. Аня дежурила 15 раз, Боря — 14 раз, Вика — 18 раз. Сколько раз дежурил Гена?

\newpage

9. Даны три числа a, b, c. Известно, что среднее арифметическое чисел a и b на 5 больше среднего арифметического всех трёх чисел. А среднее арифметическое чисел a и c на 8 меньше среднего арифметического всех трёх чисел. На сколько среднее арифметическое чисел b и c отличается от среднего арифметического всех трёх чисел?

10. Саша поехал в гости к бабушке. В субботу он сел в поезд, а через 50 часов в понедельник доехал до бабушкиного города. Саша заметил, что в этот понедельник число совпало с номером вагона, в котором он ехал, что номер его места в вагоне был меньше номера вагона и что в ту субботу, когда он садился в поезд, число было больше номера вагона. Какими были номера вагона и места? Обязательно объясните свой ответ.

11. У Винни-Пуха пять друзей, у каждого из кото- рых в домике есть горшочки с мёдом: у Тигры — 1, у Пятачка — 2, у Совы — 3, у Иа-Иа — 4, у Кролика — 5. Винни-Пух по очереди приходит в гости к каждому другу, съедает один горшо- чек мёда, а остальные забирает с собой. К последнему домику он подошёл, неся 10 горшочков с мёдом. Чей домик Пух мог посетить первым?